\chapter{Conclusions and Next Steps}

heapx is a compact but credible experimental heap library.
Its core design, one stable public API over multiple private backends, makes it possible to compare heap implementations without changing client code.
The project is intentionally narrow, but within that narrow scope it already contains many professional engineering elements.

\section{Current Strengths}

The main strengths are:

\begin{itemize}
    \item clear separation between public API and backend implementation;
    \item explicit caller ownership of stored items;
    \item generational handles for safe targeted operations;
    \item documented complexity bounds;
    \item binary, Fibonacci, and Kaplan heap backends;
    \item backend-specific invariant validation;
    \item randomized oracle tests;
    \item sanitizer and leak-check workflows;
    \item TSV benchmark output and comparison tooling;
    \item Doxygen API documentation;
    \item a LaTeX report and presentation for project communication.
\end{itemize}

\section{Most Important Next Steps}

The most valuable next steps are:

\begin{itemize}
    \item run larger benchmark campaigns on fixed DIMACS datasets;
    \item store selected benchmark summaries in a reproducible form;
    \item add optional CI jobs for additional compilers and analyzers;
    \item keep backend invariants documented as implementations evolve;
    \item evaluate whether future heap variants require API changes;
    \item avoid adding unrelated data structures that would dilute the project focus.
\end{itemize}

\section{Final Assessment}

The project is best understood as an algorithm-engineering workspace rather than only a library.
The public API provides a stable surface.
The backends expose different algorithmic trade-offs.
The tests and benchmarks make those trade-offs inspectable.

For a professional experimental repository, this is the right direction.
Future work should deepen the evidence around the existing heap implementations: correctness under stress, behavior under allocation pressure, benchmark performance across workload families, and careful documentation of any deviation between theoretical bounds and practical behavior.
